\begin{table}[!htbp]\small\setlength{\tabcolsep}{2pt}\caption*{Worked sodium bicarbonate row from Peiper and Pitzer, Table 6(a) \cite{peiperThermodynamicsAqueousCarbonate1982}.}
\begin{tabular}{@{}>{\raggedright\arraybackslash}p{0.11\linewidth}>{\raggedright\arraybackslash}p{0.1\linewidth}>{\raggedright\arraybackslash}p{0.1\linewidth}>{\raggedright\arraybackslash}p{0.1\linewidth}>{\raggedright\arraybackslash}p{0.1\linewidth}>{\raggedright\arraybackslash}p{0.1\linewidth}>{\raggedright\arraybackslash}p{0.1\linewidth}>{\raggedright\arraybackslash}p{0.1\linewidth}@{}}
\toprule $T$ (K) & $m$ & $X_2$ & $Y$ & $H_2X$ & $Z$ & $a_w$ & $\phi_{\rm eq}$ \\ \midrule
278.15 & 0.010 & 0.0001 & 0.0098 & 0.0001 & 0 & 0.99965 & 0.966 \\
\bottomrule\end{tabular}
\par\smallskip{\footnotesize\raggedright Molalities are in mol/kg water. The symbols denote nominal salt and sodium ($m$), $X_2$ carbonate, $Y$ bicarbonate, $H_2X$ dissolved \COtwo, and $Z$ hydroxide. Charge closure gives $m_{\mathrm{CO}_3^{2-}}=(m-Y-Z)/2=0.0001$. The ionic target is $[\phi_{\rm eq}(m+X_2+Y+H_2X+Z)-H_2X]/(m+m_{\mathrm{CO}_3^{2-}}+Y+Z)=(0.01932-0.00010)/0.01990=0.9658291$.\par}\end{table}
